\section{Introduction}

Artificial intelligence (AI) is increasingly embedded in production, logistics, decision support, and service delivery. As with prior general-purpose technologies, the central macroeconomic question is not whether AI can improve isolated tasks, but whether such improvements scale into sustained productivity and growth at the aggregate level \cite{bresnahanTrajtenberg1995,cominMestieri2018}. Existing policy debates frequently assume large macro gains, yet empirical evidence remains uneven because measurement of AI intensity differs across settings, adoption is heterogeneous, and national outcomes reflect many jointly evolving structural factors \cite{autor2015,brynjolfssonRockSyverson2019,imf2024,goldfarbTucker2019}.

This paper examines whether changes in AI adoption within countries are associated with changes in macroeconomic performance over time. The objective is deliberately narrow: to establish whether a stable empirical signal is visible in cross-country panel data under transparent assumptions.

\subsection{Research Questions}

\textbf{RQ1:} Is higher AI adoption associated with higher total factor productivity (TFP) within countries over time?\\
\textbf{RQ2:} Does this association also appear in near-term GDP growth dynamics?

To answer these questions, we estimate panel models that prioritize within-country variation and then evaluate whether findings persist under alternative specification and inference choices.

The contribution is threefold. First, the paper provides a fully reproducible workflow in which data validation, estimation, and manuscript outputs are linked end-to-end. Second, it reports a compact but policy-relevant robustness portfolio, including two-way fixed effects, trimmed samples, lag structures, alternative clustering, and placebo outcomes. Third, it adopts an explicit inferential boundary: estimates are interpreted as conditional associations unless stronger identification conditions are met.

A practical gap in current evidence is the mismatch between rapid AI discourse and the slower accumulation of comparable cross-country measurements. Much of the strongest evidence still comes from firm- or task-level settings with narrower domains and shorter horizons. Those studies are informative about mechanisms, but they do not map one-for-one into country-level macro elasticities. The present analysis addresses this gap by combining harmonized country-year evidence with explicit inferential boundaries, so readers can distinguish robust regularities from unresolved uncertainty.

The timing dimension is central. If AI effects diffuse through complementary assets and organizational adaptation, productivity should react before broad GDP-growth measures fully adjust. Stating this timing logic in advance improves interpretation discipline: the paper is not testing whether AI matters at all, but where in the macro transmission chain the strongest and most stable signal appears.

The manuscript is intentionally scoped as a baseline empirical contribution. It prioritizes interpretable estimates and practical robustness checks, while leaving diagnostics-heavy stress testing to the companion technical paper.

Previewing the results, we find a stable positive association between AI adoption and TFP across baseline and robustness models. The corresponding growth association is positive but smaller, consistent with the possibility that productivity effects materialize before broader macro growth responses. Sensitivity checks reinforce this pattern, while leaving open standard concerns about unobserved time-varying confounding and measurement error.

The paper is organized as follows. Section 2 situates the analysis in the technology-productivity and AI macro literatures. Section 3 describes data construction and validation. Section 4 presents econometric specifications and identifying assumptions. Section 5 reports results and model diagnostics. Section 6 discusses interpretation, limitations, and policy relevance. Section 7 concludes.
